\chapter{Barisan dan Deret Fungsi}\label{ch:b2:funcseq}

Ketika fungsi konvergen ke sebuah fungsi, sifat mana yang selamat dari
pelewatan ke limitnya? \omterm{def:b2:funcseq:def}{Kekonvergenan titik demi titik} nyaris tidak
memelihara apa pun; sedangkan kekonvergenan \emph{\omterm{def:b2:funcseq:def}{seragam}} --- yakni
kekonvergenan dalam \omterm{def:b2:nvs:norm}{norma} supremum --- memelihara kekontinuan, integral
pada ruas, dan, dengan satu pelintir, turunannya. Bab ini membuktikan
ketiga teorema pemindahannya, versi deretnya, lalu memahkotainya dengan
teorema hampiran Weierstrass, yang dibuktikan lewat \omterm{thm:b2:funcseq:weierstrass}{polinomial Bernstein}
yang indah dan bercorak peluang.

\section{Kekonvergenan titik demi titik dan seragam}

\begin{definition}\label{def:b2:funcseq:def}
Misalkan $f_n, f \colon X \to \R$ (atau $\C$, atau sebuah ruang
bernorma), dengan $X$ himpunan apa pun. Barisan $(f_n)$ konvergen ke $f$
\emph{titik demi titik}\index{kekonvergenan titik demi titik}
apabila $f_n(x) \to f(x)$ untuk setiap $x$; dan konvergen
\emph{seragam}\index{kekonvergenan seragam} apabila
\[
\norm{f_n - f}_\infty = \sup_{x \in X}\, \abs{f_n(x) - f(x)}
\xrightarrow[n \to \infty]{} 0 .
\]
Seragam mengakibatkan titik demi titik; sedangkan pada
$C(\intcc{a}{b})$, kekonvergenan seragam persis sama dengan kekonvergenan
di \omterm{def:b2:nvs:banach}{ruang Banach} $\bigl(C(\intcc{a}{b}), \norm\cdot_\infty\bigr)$ pada
\cref{ch:b2:nvs}.
\end{definition}

\begin{example}\label{ex:b2:funcseq:xn}
Pada $\intcc{0}{1}$, $f_n(x) = x^n$ konvergen \omterm{def:b2:funcseq:def}{titik demi titik} ke limit
yang \emph{tak \omterm{def:b2:metric:continuity}{kontinu}} $f = \mathbf{1}_{\{1\}}$; kekonvergenannya
tidak \omterm{def:b2:funcseq:def}{seragam}: sebab $\norm{f_n - f}_\infty \geq f_n\bigl(1 -
\tfrac1n\bigr) = (1 - \tfrac1n)^n \to \eu^{-1} \neq 0$. Namun pada
$\intcc{0}{a}$ dengan $a < 1$ ia \emph{memang} \omterm{def:b2:funcseq:def}{seragam} (sebab $\sup = a^n
\to 0$): jadi keseragaman adalah sifat daerahnya, sama seperti sifat barisannya.
\end{example}

\begin{omfigure}
\begin{tikzpicture}
\begin{axis}[omaxis, width=8.2cm, height=5.6cm,
    xmin=0, xmax=1.08, ymin=0, ymax=1.15,
    xtick={0,1}, ytick={1}]
  \addplot[omDef!40, thick, domain=0:1, samples=80] {x^2};
  \addplot[omDef!60, thick, domain=0:1, samples=80] {x^5};
  \addplot[omDef, thick, domain=0:1, samples=120] {x^15};
  \addplot[omThm, very thick, domain=0:1, samples=160] {x^40};
  \node[font=\small, omDef] at (axis cs:0.62,0.5) {$x^2$};
  \node[font=\small, omDef] at (axis cs:0.83,0.48) {$x^5$};
  \node[font=\small, omThm] at (axis cs:0.987,0.5) {$x^{40}$};
\end{axis}
\end{tikzpicture}

{\small Barisan $x^n$ pada $\intcc{0}{1}$: grafiknya melorot ke
$0$ tetapi semuanya wajib memanjat ke $1$ di $x = 1$ --- jadi jarak
supremumnya ke limit \omterm{def:b2:funcseq:def}{titik demi titik} yang tak \omterm{def:b2:metric:continuity}{kontinu} itu tak pernah
menyusut di bawah sebuah konstanta.}
\end{omfigure}

\begin{example}[Dua limit yang menolak bertukar]\label{ex:b2:funcseq:doublelimit}
Seluruh bab ini tentang menukar limit, jadi inilah kegagalan
sekecil mungkin. Misalkan $a_{n,m} = \dfrac{n}{n+m}$ untuk
$n, m \geq 1$. Maka
\[
\lim_{m\to\infty}\Bigl(\lim_{n\to\infty}a_{n,m}\Bigr)
= \lim_{m\to\infty} 1 = 1,
\qquad
\lim_{n\to\infty}\Bigl(\lim_{m\to\infty}a_{n,m}\Bigr)
= \lim_{n\to\infty} 0 = 0 :
\]
kedua limit beriterasinya ada dan keduanya berbeda. Setiap teorema
pemindahan pada bab ini adalah izin untuk menukar dua limit ---
$\lim_n$ dengan $\lim_{x\to a}$ (kekontinuan), dengan $\int$
(pengintegralan), dengan $\frac{\dd}{\dd x}$ (penurunan) ---
dan \omterm{def:b2:funcseq:def}{kekonvergenan seragam} persis menjadi bayaran yang membuat
penukaran itu sah. Pelajaran penutupnya: setiap kali sebuah ``bukti''
diam-diam menukar dua operasi limit, maka larik dua baris inilah
contoh penyangkal yang perlu diacungkan; sedangkan bonggol meluncur pada
\cref{exo:b2:funcseq:2} adalah gejala yang sama dengan busana
tanda integral.
\end{example}

\section{Ketiga teorema pemindahan}

\begin{theorem}[Kekontinuan]\label{thm:b2:funcseq:continuity}
Jika tiap $f_n$ \omterm{def:b2:metric:continuity}{kontinu} di $a$ dan $f_n \to f$ secara \omterm{def:b2:funcseq:def}{seragam} pada
sebuah persekitaran $a$, maka $f$ \omterm{def:b2:metric:continuity}{kontinu} di $a$. Jadi limit \omterm{def:b2:funcseq:def}{seragam}
fungsi \omterm{def:b2:metric:continuity}{kontinu} bersifat \omterm{def:b2:metric:continuity}{kontinu}.
\end{theorem}

\begin{proof}
Hujah $3\varepsilon$ yang sudah dipakai pada
\cref{thm:b2:metric:rncomplete}: pilih $n$ dengan $\norm{f_n -
f}_\infty \leq \varepsilon$, lalu $\delta$ dari kekontinuan
$f_n$ di $a$; maka untuk $\abs{x - a} \leq \delta$,
\[
\abs{f(x) - f(a)} \leq \abs{f(x) - f_n(x)} + \abs{f_n(x) - f_n(a)}
+ \abs{f_n(a) - f(a)} \leq 3\varepsilon . \qedhere
\]
\end{proof}

\begin{example}[Keseragaman gagal persis di tempat limitnya patah]\label{ex:b2:funcseq:threshold}
Pada $\intcc{0}{2}$, misalkan $f_n(x) = \dfrac{x^n}{1 + x^n}$. Limit
\omterm{def:b2:funcseq:def}{titik demi titiknya} adalah fungsi tiga potong:
\[
f(x) = \begin{cases} 0 & 0 \leq x < 1,\\[2pt]
\tfrac12 & x = 1,\\[2pt]
1 & 1 < x \leq 2, \end{cases}
\]
yang tak \omterm{def:b2:metric:continuity}{kontinu} di $1$, jadi menurut \cref{thm:b2:funcseq:continuity}
kekonvergenannya tak mungkin seragam pada $\intcc{0}{2}$. Namun pada
potongan tertutup yang menghindari ambangnya ia \omterm{def:b2:funcseq:def}{seragam}: untuk $0 \leq x
\leq a < 1$,
\[
\sup_{\intcc{0}{a}}\abs{f_n - 0}
= \frac{a^n}{1 + a^n} \leq a^n \to 0 ,
\]
sedangkan untuk $1 < b \leq x \leq 2$,
\[
\sup_{\intcc{b}{2}}\abs{f_n - 1}
= \frac{1}{1 + b^n} \leq b^{-n} \to 0 ,
\]
dengan kedua supremumnya terhitung lewat kemonotonan $u \mapsto
\frac{u}{1+u}$ dan $x \mapsto x^n$. Pelajaran penutupnya: kegagalan
keseragamannya terpusat di ketakkontinuan limitnya --- yakni geometri
yang sama seperti \cref{ex:b2:funcseq:xn}, sekaligus alasan mengapa
disiplin ``\omterm{def:b2:funcseq:def}{seragam} pada tiap ruas di dalamnya'' berulang sepanjang bab ini.
\end{example}

\begin{theorem}[Pengintegralan pada sebuah ruas]\label{thm:b2:funcseq:integration}
Jika $f_n \to f$ secara \omterm{def:b2:funcseq:def}{seragam} pada $\intcc{a}{b}$, dengan $f_n$
\omterm{def:b2:metric:continuity}{kontinu} sepotong-sepotong (demikian pula $f$), maka
\[
\int_a^b f_n \longrightarrow \int_a^b f .
\]
\end{theorem}

\begin{proof}
Kelinearan dan ketaksamaan segitiga untuk integral memberi
\[
\Bigl|\int_a^b f_n - \int_a^b f\Bigr|
= \Bigl|\int_a^b (f_n - f)\Bigr|
\leq \int_a^b\abs{f_n - f}
\leq (b - a)\,\norm{f_n - f}_\infty \longrightarrow 0 .
\]
Faktor panjang $(b - a)$ itulah tempat kekompakan ruasnya
masuk: sebab pada interval yang tidak \omterm{def:b2:metric:compact}{kompak}, taksiran yang sama
menghasilkan batas $\infty\cdot0$ yang tak berguna, dan kesimpulannya
sungguh gagal tanpa dominasi --- bonggol datar $f_n =
\frac1n\mathbf 1_{\intcc{0}{n}}$ konvergen \omterm{def:b2:funcseq:def}{seragam} ke $0$ pada
$\intco{0}{\infty}$ namun mempertahankan $\int f_n = 1$ (lihat catatan
jebakan di bawah), dan bonggol meluncur pada catatan
\cref{ch:b2:integration} berbuat sama dengan \omterm{def:b2:funcseq:def}{kekonvergenan titik demi
titik}; jadi \omterm{def:b2:funcseq:def}{kekonvergenan seragam} mengendalikan tinggi, tak pernah lebar.
\end{proof}

\begin{omfigure}
\begin{tikzpicture}
\begin{axis}[omaxis, width=8.6cm, height=5.4cm,
    xmin=0, xmax=1.05, ymin=0, ymax=1.9,
    xtick={0,1}, ytick={1}]
  \addplot[omDef!45, thick, domain=0:1.05, samples=120]
    {2*x*exp(-2*x^2)};
  \addplot[omDef!70, thick, domain=0:1.05, samples=120]
    {5*x*exp(-5*x^2)};
  \addplot[omDef, thick, domain=0:1.05, samples=160]
    {15*x*exp(-15*x^2)};
  \node[font=\small, omDef] at (axis cs:0.62,0.42) {$n=2$};
  \node[font=\small, omDef] at (axis cs:0.40,0.80) {$n=5$};
  \node[font=\small, omDef] at (axis cs:0.24,1.55) {$n=15$};
\end{axis}
\end{tikzpicture}

{\small Bonggol $g_n(x) = nx\,\eu^{-nx^2}$ pada
\cref{exo:b2:funcseq:1}: mereka konvergen ke $0$ di setiap titik,
tetapi puncaknya (setinggi $\sim\sqrt{n/(2\eu)}$, yang hanyut ke
$0$) tumbuh tanpa batas --- yakni \omterm{def:b2:funcseq:def}{kekonvergenan titik demi titik} dengan
$\norm{g_n}_\infty \to \infty$, dan $\int_0^1 g_n \to \frac12
\neq 0$: massanya bersembunyi di bawah puncak yang bergerak.}
\end{omfigure}

\begin{theorem}[Penurunan]\label{thm:b2:funcseq:differentiation}
Misalkan $f_n$ berkelas $C^1$ pada sebuah interval $I$, dengan: $(f_n')$
konvergen \emph{\omterm{def:b2:funcseq:def}{seragam}} pada $I$ (atau pada tiap ruas $I$) ke suatu
$g$, dan $(f_n(x_0))$ konvergen di satu titik $x_0$. Maka $(f_n)$
konvergen (\omterm{def:b2:funcseq:def}{seragam} pada tiap ruas) ke sebuah fungsi berkelas $C^1$,
sebut $f$, dan $f' = g$: jadi limitnya boleh diturunkan.
\end{theorem}

\begin{proof}
Tetapkan $f(x) = \lim f_n(x_0) + \int_{x_0}^x g$: ini sah, sebab $g$
\omterm{def:b2:metric:continuity}{kontinu} --- memang $g$ adalah limit \emph{\omterm{def:b2:funcseq:def}{seragam}} pada ruas atas $f_n'$
yang \omterm{def:b2:metric:continuity}{kontinu}, jadi \cref{thm:b2:funcseq:continuity} berlaku, dan
integral sebuah fungsi \omterm{def:b2:metric:continuity}{kontinu} terdefinisi dengan baik dengan, menurut
teorema dasar kalkulus,
\[
f'(x) = g(x) \qquad (x \in I) :
\]
jadi calon limitnya berkelas $C^1$ dengan turunan yang benar
\emph{menurut konstruksinya}, sebelum kekonvergenan apa pun dibuktikan.
Menurut teorema dasar itu lagi, $f_n(x) = f_n(x_0) + \int_{x_0}^x
f_n'$; setelah dikurangkan,
\[
\abs{f_n(x) - f(x)} \leq \abs{f_n(x_0) - \lim f_n(x_0)}
+ \abs{x - x_0}\,\norm{f_n' - g}_{\infty} ,
\]
yang menuju $0$ secara \omterm{def:b2:funcseq:def}{seragam} pada tiap ruas. Dan $f$ berkelas $C^1$
dengan $f' = g$ menurut konstruksinya.
\end{proof}

\begin{example}[Mengapa hipotesisnya duduk pada turunannya]\label{ex:b2:funcseq:sqrtabs}
Misalkan $F_n(x) = \sqrt{x^2 + \frac1n}$ pada $\R$. Tiap $F_n$ berkelas
$C^1$ (bahkan $C^\infty$), dan kekonvergenannya ke $\abs x$ bersifat
\omterm{def:b2:funcseq:def}{seragam} pada seluruh $\R$:
\[
0 \leq F_n(x) - \abs x
= \frac{(x^2 + \frac1n) - x^2}{\sqrt{x^2+\frac1n} + \abs x}
= \frac{1/n}{\sqrt{x^2 + \frac1n} + \abs x}
\leq \frac{1/n}{1/\sqrt n} = \frac{1}{\sqrt n} .
\]
Namun limitnya $\abs x$ tak dapat diturunkan di $0$: jadi \omterm{def:b2:funcseq:def}{kekonvergenan
seragam} \emph{fungsinya}, secepat apa pun, tidak memindahkan
keterdiferensialan. Kegagalannya kasatmata pada turunannya:
\[
F_n'(x) = \frac{x}{\sqrt{x^2 + \frac1n}}
\longrightarrow \begin{cases} 1 & x > 0,\\ 0 & x = 0,\\
-1 & x < 0, \end{cases}
\]
yakni limit \omterm{def:b2:funcseq:def}{titik demi titik} yang tak \omterm{def:b2:metric:continuity}{kontinu}, sehingga $(F_n')$ tak
mungkin konvergen \omterm{def:b2:funcseq:def}{seragam} di dekat $0$ (lagi-lagi
\cref{thm:b2:funcseq:continuity}). Pelajaran penutupnya:
\cref{thm:b2:funcseq:differentiation} dengan sengaja mengandaikan
\omterm{def:b2:funcseq:def}{kekonvergenan seragam} pada $f_n'$, bukan pada $f_n$ ---
dan contoh inilah alasannya.
\end{example}

\section{Deret fungsi}

\begin{definition}\label{def:b2:funcseq:series}
Sebuah deret fungsi $\sum u_n$ konvergen \omterm{def:b2:funcseq:def}{titik demi titik} atau \omterm{def:b2:funcseq:def}{seragam}
apabila jumlah parsialnya demikian. Ia konvergen
\emph{normal}\index{kekonvergenan normal} (pada $X$) apabila $\sum
\norm{u_n}_\infty < \infty$. Kekonvergenan normal mengakibatkan
\omterm{def:b2:funcseq:def}{kekonvergenan seragam} (di \omterm{def:b2:nvs:banach}{ruang Banach} berisi fungsi terbatas:
\cref{thm:b2:nvs:absoluteconvergence}), yang mengakibatkan \omterm{def:b2:funcseq:def}{kekonvergenan
titik demi titik}; dan kedua implikasinya tegas.
\end{definition}

\begin{example}[Satu deret, tiga vonis]\label{ex:b2:funcseq:threemodes}
Ambil $u_n(x) = \frac{x^n}{n}$ pada $\intco{0}{1}$. \emph{\omterm{def:b2:funcseq:def}{Titik demi
titik}:} ia konvergen untuk setiap $x \in \intco01$ (lewat pembandingan
dengan deret geometri). \emph{\omterm{def:b2:funcseq:series}{Normal} pada $\intcc{0}{a}$ dengan $a < 1$:}
sebab $\norm{u_n}_{\infty,\intcc0a} = \frac{a^n}{n}$, yang \omterm{def:b2:series:summable}{terjumlahkan}.
\emph{Tidak \omterm{def:b2:funcseq:series}{normal} pada $\intco{0}{1}$:} sebab
$\norm{u_n}_{\infty,\intco01} = \frac1n$, dan $\sum\frac1n$
divergen. \emph{Bahkan tidak \omterm{def:b2:funcseq:def}{seragam} pada $\intco{0}{1}$:} sebab
sisanya membandel di dekat $1$,
\[
R_N(x) = \sum_{n>N}\frac{x^n}{n}
\geq \sum_{n=N+1}^{2N}\frac{x^n}{n}
\geq \frac{N\,x^{2N}}{2N} = \frac{x^{2N}}{2}
\xrightarrow[x\to1^-]{} \frac12 ,
\]
jadi $\sup_{\intco01}\abs{R_N} \geq \frac12$ untuk setiap $N$.
Pelajaran penutupnya: keempat vonis itu hidup rukun --- sebab
jumlahnya $-\ln(1-x)$ \omterm{def:b2:metric:continuity}{kontinu} pada $\intco{0}{1}$ karena
kekontinuan hanya memerlukan keseragaman \emph{di dekat tiap titik},
yakni pada ruas $\intcc0a$; sedangkan meledak di tepinya adalah
hak jumlah itu sendiri.
\end{example}

\begin{theorem}[Pemindahan untuk deret]\label{thm:b2:funcseq:seriestransfer}
Jika $\sum u_n$ konvergen \omterm{def:b2:funcseq:def}{seragam} (misalnya \omterm{def:b2:funcseq:series}{normal}) pada himpunan yang
bersangkutan: maka kekontinuan semua $u_n$ di $a$ berpindah ke jumlahnya;
pengintegralan pada sebuah ruas dapat dilakukan suku demi suku; dan jika
$\sum u_n(x_0)$ konvergen sedangkan $\sum u_n'$ konvergen \omterm{def:b2:funcseq:def}{seragam} pada
tiap ruas, maka jumlahnya berkelas $C^1$ dengan turunan $\sum u_n'$.
\end{theorem}

\begin{proof}
Semuanya tak lain teorema yang bersesuaian yang diterapkan pada jumlah
parsial $S_N = \sum_{n\leq N}u_n$, yang merupakan jumlah hingga atas
fungsi dengan keteraturan yang bersangkutan. Kekontinuan: tiap $S_N$
\omterm{def:b2:metric:continuity}{kontinu} di $a$ dan $S_N \to \sum u_n$ secara \omterm{def:b2:funcseq:def}{seragam}, jadi
\cref{thm:b2:funcseq:continuity}. Pengintegralan: pada ruas itu,
\[
\int_a^b \sum_{n\geq0} u_n
= \lim_N \int_a^b S_N
= \lim_N \sum_{n=0}^{N}\int_a^b u_n
= \sum_{n\geq0}\int_a^b u_n
\]
menurut \cref{thm:b2:funcseq:integration} (kesamaan pertamanya) dan
kelinearan integralnya (kesamaan keduanya). Penurunan: tiap
$S_N$ berkelas $C^1$, $S_N(x_0)$ konvergen, dan $S_N' =
\sum_{n\leq N}u_n'$ konvergen \omterm{def:b2:funcseq:def}{seragam} pada tiap ruas, jadi
\cref{thm:b2:funcseq:differentiation} memberi bahwa jumlahnya berkelas
$C^1$ dengan turunan $\lim S_N' = \sum u_n'$.
\end{proof}

\begin{remark}[Jebakan yang sering muncul]
Ada empat jebakan, semuanya terlihat di lembar ujian. \emph{(i)
Supremum yang diperiksa separuh:} menilai $f_n$ sepanjang barisan
$x_n$ yang dipilih baik-baik hanya membatasi $\norm{f_n - f}_\infty$
dari \emph{bawah} --- cukup untuk menyangkal keseragamannya (seperti pada
\cref{ex:b2:funcseq:xn}), tak pernah untuk membuktikannya; dan untuk
membuktikannya, batasilah supremumnya lewat perhitungan yang sahih untuk
\emph{setiap} $x$. \emph{(ii) Keseragaman pada himpunan yang salah:}
\omterm{def:b2:funcseq:series}{kekonvergenan normal} atau \omterm{def:b2:funcseq:def}{seragam} sering berlaku pada tiap
$\intcc{-a}{a}$ atau $\intco\delta\infty$ tetapi gagal pada gabungan
\omterm{def:b2:metric:topology}{terbukanya}; hal itu bukan halangan --- sebab kekontinuan dan
keterdiferensialan bersifat lokal, jadi disiplin ruas demi ruas pada
\cref{ex:b2:funcseq:zeta} memberikannya pada seluruh himpunan \omterm{def:b2:metric:topology}{terbuka}
itu. \emph{(iii) Mengintegralkan pada yang bukan ruas:}
\cref{thm:b2:funcseq:integration} adalah pernyataan tentang
ruas; sedangkan pada $\intco0\infty$, \omterm{def:b2:funcseq:def}{kekonvergenan seragam} tidak
mencegah massanya lolos ke tak hingga (sebab $f_n =
\frac1n\mathbf 1_{\intcc{0}{n}}$ konvergen \omterm{def:b2:funcseq:def}{seragam} ke $0$
dengan $\int f_n = 1$) --- jadi pakailah kekonvergenan terdominasi di
sana. \emph{(iv) Menurunkan limitnya:}
\cref{ex:b2:funcseq:sqrtabs}; hipotesis turunannya ada pada
$(f_n')$, dan laju kekonvergenan $(f_n)$ mana pun tak dapat
menggantikannya.
\end{remark}

\begin{example}[Fungsi $\zeta$ Riemann]\label{ex:b2:funcseq:zeta}
Deret $\zeta(s) = \sum_{n\geq1} n^{-s}$ konvergen \omterm{def:b2:funcseq:series}{normal} pada tiap
setengah-garis $\intco{a}{+\infty}$ dengan $a > 1$ (sebab
$\norm{n^{-s}}_\infty = n^{-a}$, yang \omterm{def:b2:series:summable}{terjumlahkan}): jadi $\zeta$
\omterm{def:b2:metric:continuity}{kontinu} pada $\intoo{1}{+\infty}$; lalu setelah diturunkan suku demi
suku (sebab deret turunannya $\sum -\ln n\; n^{-s}$ juga konvergen
\omterm{def:b2:funcseq:series}{normal} pada $\intco{a}{\infty}$), $\zeta$ berkelas $C^1$ --- dan, setelah
diiterasikan, $C^\infty$ --- dengan $\zeta'(s) = -\sum \frac{\ln
n}{n^s}$. Perhatikan disiplinnya: \omterm{def:b2:funcseq:series}{kekonvergenan normalnya} diperiksa pada
\emph{sub}-setengah-garis, tak pernah pada $\intoo{1}{\infty}$ yang
\omterm{def:b2:metric:topology}{terbuka}, tempat ia gagal.
\end{example}

\begin{example}[Sebuah deret logaritmik, dikerjakan sampai tuntas]\label{ex:b2:funcseq:logseries}
Misalkan $F(x) = \sum_{n\geq1} \frac{\eu^{-nx}}{n}$ pada
$\intoo{0}{\infty}$. Tiap sukunya terbatas pada $\intco{\delta}\infty$
oleh $\frac{\eu^{-n\delta}}{n} \leq \eu^{-n\delta}$, yakni deret
geometri yang konvergen: jadi kekonvergenannya \omterm{def:b2:funcseq:series}{normal} pada tiap
$\intco\delta\infty$, sehingga $F$ \omterm{def:b2:metric:continuity}{kontinu} pada
$\intoo{0}{\infty}$. Deret turunannya $\sum -\eu^{-nx}$ juga
konvergen \omterm{def:b2:funcseq:series}{normal} pada $\intco\delta\infty$ (sebab
$\norm{\eu^{-nx}}_{\infty,\intco\delta\infty} = \eu^{-n\delta}$), jadi
$F$ berkelas $C^1$ dengan turunan yang geometri:
\[
F'(x) = -\sum_{n\geq1}\eu^{-nx}
= \frac{-\eu^{-x}}{1 - \eu^{-x}}
= \frac{-1}{\eu^{x} - 1} .
\]
Setelah diiterasikan, $F$ berkelas $C^\infty$. Dengan mengintegralkan
$F'$ (sebab baik $F$ maupun $x \mapsto -\ln(1 - \eu^{-x})$ bernilai nol
di $+\infty$ dan punya turunan yang sama pada $\intoo0\infty$):
\[
F(x) = -\ln\bigl(1 - \eu^{-x}\bigr),
\]
yakni deret logaritmik di $t = \eu^{-x}$. Pelajaran penutupnya: ketika $x
\to 0^+$, $F(x) = -\ln(x + O(x^2)) = \ln\frac1x + O(x)$ --- jadi
deretnya divergen secara logaritmik di perbatasannya, persis seperti
deret harmonik yang menjadi wujudnya di $x = 0$; dan \omterm{def:b2:funcseq:series}{kekonvergenan
normal} pada $\intco\delta\infty$ tetapi tidak pada $\intoo0\infty$
adalah gejalanya.
\end{example}

\begin{method}[Membuktikan atau menyangkal kekonvergenan seragam]\label{met:b2:funcseq:uniform}
Untuk $f_n \to f$ \omterm{def:b2:funcseq:def}{titik demi titik} pada $X$:
\begin{enumerate}
  \item Hitung atau batasi $\norm{f_n - f}_\infty$: telaahlah
        fungsi $x \mapsto \abs{f_n(x) - f(x)}$ (lewat turunan atau
        kemonotonan) untuk menemukan maksimumnya; sebab batas yang
        sahih untuk setiap $x$ dan menuju $0$ membuktikan keseragamannya.
  \item Untuk \emph{menyangkalnya}: tunjukkan titik $x_n$ dengan
        $\abs{f_n(x_n) - f(x_n)} \not\to 0$ (sering $x_n$ mengikuti
        bonggol yang bergerak, seperti pada \cref{exo:b2:funcseq:1});
        atau panggillah sebuah teorema pemindahan secara kontrapositif
        --- yakni limit tak \omterm{def:b2:metric:continuity}{kontinu} dari fungsi \omterm{def:b2:metric:continuity}{kontinu}
        (\cref{ex:b2:funcseq:threshold}), atau $\int f_n
        \not\to \int f$ pada sebuah ruas.
  \item Untuk deret, cobalah \omterm{def:b2:funcseq:series}{kekonvergenan normal} lebih dulu
        (yakni $\sum\sup\abs{u_n} < \infty$); jika gagal secara global,
        ujilah pada subruas yang penting
        (\cref{ex:b2:funcseq:threemodes}); dan jika gagal
        di mana-mana, \omterm{def:b2:funcseq:def}{kekonvergenan seragamnya} masih mungkin berlaku
        lewat batas sisa berselang-seling
        (\cref{exo:b2:funcseq:4}) atau lewat penjumlahan parsial.
\end{enumerate}
\end{method}

\section{Teorema hampiran Weierstrass}

\begin{theorem}[Weierstrass, lewat Bernstein]\label{thm:b2:funcseq:weierstrass}
Setiap $f \colon \intcc{0}{1} \to \R$ yang \omterm{def:b2:metric:continuity}{kontinu} merupakan limit
\omterm{def:b2:funcseq:def}{seragam} polinomial --- lebih tepatnya, limit \emph{polinomial
Bernstein}\index{polinomial Bernstein}\index{teorema hampiran
Weierstrass} miliknya sendiri
\[
B_n(f)(x) = \sum_{k=0}^{n} f\Bigl(\frac kn\Bigr)\binom nk x^k
(1-x)^{n-k} .
\]
\end{theorem}

\begin{proof}
Tetapkan $x \in \intcc{0}{1}$ lalu tetapkan $p_k(x) = \binom nk x^k(1 -
x)^{n-k}$. Ada tiga kesamaan binomial, yang diperoleh dengan menilai
$(x + y)^n$ beserta dua turunannya terhadap $x$ di $y = 1 - x$:
\[
\sum_k p_k = 1,
\qquad
\sum_k k\,p_k = nx,
\qquad
\sum_k k(k-1) p_k = n(n-1)x^2 .
\]
Secara rinci: $(x+y)^n = \sum_k\binom nk x^ky^{n-k}$ di $y = 1-x$
adalah yang pertama; lalu menurunkannya terhadap $x$,
\[
n(x+y)^{n-1} = \sum_k k\binom nk x^{k-1}y^{n-k} ,
\]
lalu mengalikannya dengan $x$ dan mengambil $y = 1 - x$ memberi
yang kedua; sedangkan menurunkannya dua kali lalu mengalikan dengan
$x^2$ memberi yang ketiga. Setelah $(k - nx)^2 = k(k-1) + k(1 - 2nx) +
n^2x^2$ diuraikan lalu ketiganya dipadukan:
\[
\sum_k (k - nx)^2 p_k
= n(n-1)x^2 + nx(1 - 2nx) + n^2x^2
= nx(1 - x) \leq \frac n4 ,
\]
yakni \emph{kesamaan ragam}.

Kini taksirlah, dengan memakai $\sum p_k = 1$:
\[
\abs{B_n(f)(x) - f(x)}
\leq \sum_{k} \Bigl| f\Bigl(\frac kn\Bigr) - f(x)\Bigr|\, p_k(x)
= \Sigma_{\text{dekat}} + \Sigma_{\text{jauh}} ,
\]
dengan pemecahannya menurut $\abs{\frac kn - x} \leq \delta$ atau tidak.
Diberikan $\varepsilon > 0$, kekontinuan seragam $f$ (lewat Heine)
menyediakan $\delta$ dengan $\Sigma_{\text{dekat}} \leq \varepsilon$.
Untuk jumlah yang jauh, dengan $M = \norm f_\infty$: menurut kesamaan
ragam dan siasat pencacahan Chebyshev,
\[
\Sigma_{\text{jauh}} \leq 2M \sum_{\abs{k - nx} > n\delta} p_k
\leq 2M\,\frac{\sum_k (k - nx)^2 p_k}{n^2\delta^2}
\leq \frac{2M}{4 n \delta^2} \cdot 1
= \frac{M}{2n\delta^2}
\xrightarrow[n\to\infty]{} 0 ,
\]
secara \omterm{def:b2:funcseq:def}{seragam} pada $x$. Jadi $\norm{B_n(f) - f}_\infty \leq \varepsilon +
\frac{M}{2n\delta^2} \leq 2\varepsilon$ untuk $n$ yang besar.
\end{proof}

\begin{remark}
Lewat penyulihan afin, teoremanya berlaku pada ruas $\intcc{a}{b}$ mana
pun. Ia gagal pada $\R$ (sebab limit \omterm{def:b2:funcseq:def}{seragam} polinomial pada
$\R$ tetaplah polinomial: \cref{exo:b2:funcseq:8}). Adapun pembacaan
bercorak peluangnya --- yakni $B_n(f)(x)$ merupakan nilai harapan $f$ di
sebuah rata-rata binomial, dan batas ragamnya adalah ketaksamaan
Chebyshev --- dijujurkan pada \cref{ch:b2:genfun}.
\end{remark}

\begin{omfigure}
\begin{tikzpicture}
\begin{axis}[omaxis, width=8.6cm, height=6.0cm,
    xmin=0, xmax=1.05, ymin=0, ymax=1.08,
    xtick={0,1}, ytick={1}]
  \addplot[omThm, very thick, domain=0:1, samples=80] {x^2};
  \addplot[omDef!45, thick, domain=0:1, samples=2] {x};
  \addplot[omDef!70, thick, domain=0:1, samples=80]
    {x^2 + x*(1-x)/2};
  \addplot[omDef, thick, domain=0:1, samples=80]
    {x^2 + x*(1-x)/6};
  \node[font=\small, omDef] at (axis cs:0.30,0.37) {$B_1f$};
  \node[font=\small, omDef] at (axis cs:0.52,0.46) {$B_2f$};
  \node[font=\small, omDef] at (axis cs:0.76,0.56) {$B_6f$};
  \node[font=\small, omThm] at (axis cs:0.88,0.68) {$f$};
\end{axis}
\end{tikzpicture}

{\small Hampiran Bernstein bagi $f(x) = x^2$ (merah), dengan memakai
rumus eksak $B_nf = x^2 + \frac{x(1-x)}{n}$ pada
\cref{exo:b2:funcseq:7}: $B_1f$ adalah tali busurnya, dan tiap kali $n$
dilipatduakan jurangnya terparuh. Andal tetapi lamban --- yakni
penjenuhan $\frac1n$ yang dieksakkan teorema Voronovskaya (soal akhir pekan).}
\end{omfigure}

\begin{remark}[Di mana ini dipakai]
Hampiran Weierstrass adalah teorema kepadatan bagi analisis
klasik: ia membuat $C(\intcc ab)$ terpisahkan, membiarkan kita memeriksa
kesamaan integral pada polinomial saja (yakni masalah momen), dan
melandasi versi trigonometrinya yang dibuktikan pada bab Fourier lewat
kernel Fejér. Soal akhir pekan bab ini menyarikan
isi kuantitatif bukti Bernstein --- yakni laju kekonvergenan yang
dikuasai modulus kekontinuannya --- lalu memisahkan apa yang sebenarnya
membuatnya berhasil, di dalam teorema Korovkin: yakni kepositifan
ditambah tiga fungsi uji. Jilid Tahun ke-3 menyeluruhkan pernyataan
kepadatannya ke subaljabar sembarang (Stone--Weierstrass) dan ke ruang \omterm{def:b2:metric:compact}{kompak}.
\end{remark}

\begin{example}[Hampiran poligonal, beserta lajunya]\label{ex:b2:funcseq:polygonal}
Untuk $L$ sebagai konstanta \omterm{def:b2:metric:continuity}{Lipschitz} bagi $f$ pada $\intcc{0}{1}$, misalkan $I_nf$
penginterpolasi afin sepotong-sepotong pada simpul $\frac kn$. Pada
sebuah sel $\intcc{\frac kn}{\frac{k+1}n}$, baik $f(x)$ maupun
$I_nf(x)$ berada di antara nilai ekstrem yang dapat diambil fungsi
Lipschitz-$L$ dengan kedua nilai simpulnya, jadi untuk $x$ di
sel itu, dengan menulis $x_k = \frac kn$:
\[
\abs{I_nf(x) - f(x)}
\leq \abs{I_nf(x) - f(x_k)} + \abs{f(x_k) - f(x)}
\leq L\,\abs{x - x_k} + L\,\abs{x - x_k}
\leq \frac{2L}{n}
\]
(sebab penginterpolasinya sendiri Lipschitz-$L$ pada sel itu: sebab
kemiringannya adalah hasil bagi selisih $f$). Karenanya $\norm{I_nf -
f}_\infty \leq \frac{2L}{n}$: jadi hampiran poligonal bagi fungsi
\omterm{def:b2:metric:continuity}{Lipschitz} konvergen dengan laju $\frac1n$ ---
\emph{lebih cepat} daripada $\frac{1}{\sqrt n}$ milik Bernstein untuk
kelas yang sama (soal akhir pekan, Bagian II). Pelajaran penutupnya:
poligonnya menginterpolasi tetapi tidak mulus, Bernstein mulus
tetapi lamban; jadi tak ada makan siang gratis antara keteraturan
penghampirnya dan lajunya --- yakni pertukaran yang dicermatkan
hasil penjenuhan pada soal akhir pekan.
\end{example}

\begin{remark}[Pandangan ke depan di dalam jilid ini]
\omterm{def:b2:funcseq:def}{Kekonvergenan seragam} menjadi kuda beban buku ini mulai sekarang. Bab
deret pangkat seluruhnya berjalan pada \omterm{def:b2:funcseq:series}{kekonvergenan normal} pada
subcakram \omterm{def:b2:metric:compact}{kompak} --- setiap teorema suku demi suku di sana adalah
kasus khusus teorema pemindahan bab ini. Bab Fourier
hidup satu lantai di atasnya: jumlah parsialnya $S_N$ gagal
persis di tempat bab ini memperingatkan (yakni \omterm{def:b2:funcseq:def}{titik demi titik} tetapi
tidak \omterm{def:b2:funcseq:def}{seragam} di lompatannya), dan rata-rata Fejér-nya berhasil lewat
mekanika $3\varepsilon$ yang sama seperti yang membuktikan
\cref{thm:b2:funcseq:continuity}. Bab persamaan diferensial
menetapkan $\eu^{tA}$ lewat deret yang konvergen \omterm{def:b2:funcseq:series}{normal} lalu
menurunkannya suku demi suku --- yakni harfiah
\cref{thm:b2:funcseq:seriestransfer} yang diterapkan pada unsur
matriksnya. Bila kelak muncul keraguan di dalam buku ini tentang
``mengapa kita boleh melakukan ini'', jawabannya biasanya sebuah
teorema bab ini.
\end{remark}

\section{Latihan}

\begin{exercise}[$\star$]\label{exo:b2:funcseq:1}
Telaahlah \omterm{def:b2:funcseq:def}{kekonvergenan titik demi titik} dan \omterm{def:b2:funcseq:def}{kekonvergenan seragamnya}
pada $\intcc{0}{1}$, lalu pada $\intcc{0}{a}$ ($a < 1$) atau
$\intco{\delta}{1}$ sesuai keperluan, bagi:
\[
f_n(x) = \frac{x}{1 + nx},
\qquad
g_n(x) = n x\,\eu^{-n x^2},
\qquad
h_n(x) = x^n(1 - x^n).
\]
\end{exercise}

\begin{exercise}[$\star$]\label{exo:b2:funcseq:2}
Buktikan bahwa $\displaystyle\int_0^1 g_n \not\to \int_0^1 \lim g_n$
bagi $g_n$ pada \cref{exo:b2:funcseq:1}, lalu damaikan dengan
\cref{thm:b2:funcseq:integration}.
\end{exercise}

\begin{exercise}[$\star$]\label{exo:b2:funcseq:3}
Buktikan bahwa $S(x) = \sum_{n\geq1} \dfrac{x^n}{n^2}$ \omterm{def:b2:metric:continuity}{kontinu} pada
$\intcc{-1}{1}$, dan bahwa $S$ berkelas $C^1$ pada $\intoo{-1}{1}$ dengan
$S'(x) = -\frac{\ln(1-x)}{x}$ untuk $0 < \abs x < 1$.
\end{exercise}

\begin{exercise}[$\star\star$]\label{exo:b2:funcseq:4}
Misalkan $F(x) = \sum_{n \geq 0} \dfrac{(-1)^n}{n + x}$ pada
$\intoo{0}{+\infty}$. Buktikan \omterm{def:b2:funcseq:def}{kekonvergenan seragamnya} (tetapi tidak
\omterm{def:b2:funcseq:series}{normal}) pada $\intco{\delta}{\infty}$ lewat batas sisa deret
berselang-seling, lalu kekontinuannya, dan persamaan fungsionalnya $F(x)
+ F(x + 1) = \frac1x$.
\end{exercise}

\begin{exercise}[$\star\star$]\label{exo:b2:funcseq:5}
(Dini) Misalkan $f_n \colon K \to \R$ \omterm{def:b2:metric:continuity}{kontinu} pada \omterm{def:b2:metric:def}{ruang metrik} yang
\emph{\omterm{def:b2:metric:compact}{kompak}}, dengan $f_n \to f$ \emph{\omterm{def:b2:funcseq:def}{titik demi titik}}, $f$ \omterm{def:b2:metric:continuity}{kontinu},
dan $(f_n(x))$ \emph{turun} terhadap $n$ untuk tiap $x$. Buktikan bahwa
kekonvergenannya \omterm{def:b2:funcseq:def}{seragam}. \emph{(Diberikan $\varepsilon$, himpunan
\omterm{def:b2:metric:topology}{terbuka} $U_n = \{x : f_n(x) - f(x) < \varepsilon\}$ naik dan menyelimuti
$K$; lalu sarikan subselimut hingganya ---
\cref{thm:b2:metric:borellebesgue}.)}
\end{exercise}

\begin{exercise}[$\star\star$]\label{exo:b2:funcseq:6}
Buktikan bahwa $\displaystyle\lim_{n\to\infty} \int_0^1
\frac{n\,f(x)}{1 + n^2x^2}\,\dd x = \frac{\pi}{2} f(0)$ untuk setiap
$f$ yang \omterm{def:b2:metric:continuity}{kontinu} pada $\intcc{0}{1}$. \emph{(Sulihkan $u = nx$; pisahkan
$f(0)$; lalu dominasilah.)}
\end{exercise}

\begin{exercise}[$\star\star$]\label{exo:b2:funcseq:7}
Hitunglah \omterm{thm:b2:funcseq:weierstrass}{polinomial Bernstein} bagi $f(x) = x^2$ secara gamblang lalu
periksalah galat \omterm{def:b2:funcseq:def}{seragamnya} $\norm{B_n f - f}_\infty =
O\bigl(\frac1n\bigr)$ yang diramalkan bukti
\cref{thm:b2:funcseq:weierstrass} --- yang di sini tepat $\frac{x(1 -
x)}{n}$ di tiap titiknya.
\end{exercise}

\begin{exercise}[$\star\star$]\label{exo:b2:funcseq:8}
Buktikan bahwa jika polinomial $P_n$ konvergen \omterm{def:b2:funcseq:def}{seragam} \emph{pada
seluruh $\R$} ke $f$, maka $f$ sebuah polinomial. \emph{(Untuk $m, n$
yang besar, $P_n - P_m$ merupakan polinomial terbatas pada $\R$, jadi
konstan; sehingga barisannya mapan kecuali konstanta.)}
\end{exercise}

\begin{exercise}[$\star\star\star$]\label{exo:b2:funcseq:9}
(Sebuah fungsi \omterm{def:b2:metric:continuity}{kontinu} yang tak dapat diturunkan di mana pun ---
berpandu) Misalkan $\varphi$ jarak ke bilangan bulat terdekat (yang
berperiode $1$, dengan $\norm{\varphi}_\infty = \frac12$, dan
Lipschitz-$1$) dan
\[
W(x) = \sum_{n=0}^{\infty} \Bigl(\frac{3}{4}\Bigr)^{\!n}
\varphi(4^n x) .
\]
Buktikan: (a) $W$ \omterm{def:b2:metric:continuity}{kontinu} pada $\R$ (lewat \omterm{def:b2:funcseq:series}{kekonvergenan normal}); (b)
untuk setiap $x$ dan setiap $m$, dengan memilih $h_m = \pm\frac12\cdot
4^{-m}$ bertanda yang membuat $\varphi$ afin pada ruas dari $4^m x$
sampai $4^m(x + h_m)$, hasil bagi selisihnya memenuhi
\[
\Bigl|\frac{W(x + h_m) - W(x)}{h_m}\Bigr| \geq 3^m -
\sum_{n<m} 3^n \geq \frac{3^m + 1}{2} \xrightarrow[m\to\infty]{}
\infty
\]
(sebab suku $n > m$ lenyap berkat keperiodikannya; suku $n = m$
menyumbang tepat $3^m$; dan suku $n < m$ terbatas berkat sifat
\omterm{def:b2:metric:continuity}{Lipschitznya}). Simpulkan bahwa $W$ tak dapat diturunkan di mana pun.
\end{exercise}

\begin{exercise}[$\star$]\label{exo:b2:funcseq:10}
Misalkan $u_n(x) = (-1)^n x^n(1 - x)$ pada $\intcc{0}{1}$. Tunjukkan
bahwa $\sum u_n$ konvergen \omterm{def:b2:funcseq:def}{titik demi titik} pada $\intcc{0}{1}$ lalu
hitung jumlahnya; tunjukkan bahwa kekonvergenannya \omterm{def:b2:funcseq:def}{seragam} pada
$\intcc{0}{1}$ \emph{(batasi sisanya $R_N(x) = \sum_{n > N} u_n(x)$,
yakni ekor geometri, oleh suku pertamanya lalu maksimumkan
$x^{N+1}(1-x)$)} tetapi \emph{tidak} \omterm{def:b2:funcseq:series}{normal} \emph{(hitung
$\norm{u_n}_\infty$)}: jadi \omterm{def:b2:funcseq:def}{kekonvergenan seragam} tegas lebih lemah
daripada \omterm{def:b2:funcseq:series}{kekonvergenan normal}. Bandingkan dengan $\sum x^n(1-x)$, yang
jumlahnya tak \omterm{def:b2:metric:continuity}{kontinu} di $1$: di sana keseragamannya pun gagal.
\end{exercise}

\begin{exercise}[$\star\star$]\label{exo:b2:funcseq:11}
Misalkan $f_n \to f$ secara \omterm{def:b2:funcseq:def}{seragam} pada sebuah \omterm{def:b2:metric:def}{ruang metrik} $X$, dengan
tiap $f_n$ \omterm{def:b2:metric:continuity}{kontinu}, dan misalkan $x_n \to x$ di $X$. Buktikan $f_n(x_n)
\to f(x)$. Tunjukkan lewat sebuah contoh pada $X = \intcc{0}{1}$ bahwa
\omterm{def:b2:funcseq:def}{kekonvergenan titik demi titik} tidak cukup, bahkan dengan $f$ yang
\omterm{def:b2:metric:continuity}{kontinu} \emph{(pakailah bonggol $g_n$ pada \cref{exo:b2:funcseq:1} dan
$x_n = \frac{1}{\sqrt{2n}}$)}.
\end{exercise}

\begin{exercise}[$\star\star\star$]\label{exo:b2:funcseq:12}
(Sebuah persamaan integral Volterra lewat deret) Untuk $f \in
C(\intcc{0}{1})$ tetapkan $Tf(x) = \int_0^x f(t)\,\dd t$.
\begin{enumerate}
  \item Tunjukkan secara induktif bahwa untuk $n \geq 1$:
        \[
        T^n f(x) = \int_0^x \frac{(x -
        t)^{n-1}}{(n-1)!}\,f(t)\,\dd t,
        \qquad
        \norm{T^n f}_\infty \leq \frac{\norm f_\infty}{n!} .
        \]
  \item Turunkan bahwa $S = \sum_{n\geq0} T^n f$ konvergen \omterm{def:b2:funcseq:series}{normal}
        pada $\intcc{0}{1}$ dan menyelesaikan persamaan integral $S = f
        + TS$.
  \item Periksalah bahwa $S(x) = f(x) + \int_0^x \eu^{x-t}f(t)\,\dd
        t$ menyelesaikan persamaan yang sama, lalu buktikan ketunggalan
        penyelesaian yang \omterm{def:b2:metric:continuity}{kontinu} \emph{(sebab jika $S = TS$ maka
        $\norm{S}_\infty \leq \norm{T^nS}_\infty \to 0$)}:
        lalu simpulkan bentuk tertutup jumlahnya.
\end{enumerate}
\end{exercise}

\section{Soal: Laju hampiran dan teorema Korovkin}

\begin{problem}\label{pb:b2:funcseq:1}
Bukti Bernstein bagi \cref{thm:b2:funcseq:weierstrass} menyembunyikan
dua harta. Pertama, ia bersifat \emph{kuantitatif}: seberapa cepat $B_nf
\to f$ dikuasai modulus kekontinuan $f$, dengan laju tajamnya
tercapai oleh $\abs{x - \frac12}$. Kedua, ia bersifat
\emph{struktural}: yang penting hanyalah bahwa $B_n$ merupakan
operator linear positif yang berkelakuan baik pada $1$, $x$, $x^2$ ---
dan pengamatan itu, setelah dipisahkan, adalah \emph{teorema
Korovkin}\index{teorema Korovkin}. Soal ini membuktikan keduanya,
lalu ditutup dengan asimtotik eksak Voronovskaya. Di sepanjang soal ini,
$f \in C(\intcc{0}{1})$, $M = \norm f_\infty$, $p_k(x) = \binom
nk x^k(1-x)^{n-k}$, dan $e_j$ menyatakan $x \mapsto x^j$.

\medskip
\noindent\textbf{Bagian I --- Operator Bernstein.}
\begin{enumerate}
  \item Tunjukkan bahwa $B_n$ bersifat linear, \emph{positif} (yakni $f
        \geq 0 \Rightarrow B_nf \geq 0$), jadi monoton ($f \leq g
        \Rightarrow B_nf \leq B_ng$), dengan
        $\norm{B_nf}_\infty \leq \norm f_\infty$, dan bahwa
        $B_nf$ menginterpolasi $f$ di kedua ujungnya.
  \item Turunkan kembali kesamaan $B_n e_0 = e_0$, $B_n e_1 =
        e_1$ dan $B_n e_2 = e_2 + \frac{e_1 - e_2}{n}$
        \emph{(turunkan $(x + y)^n$ dua kali lalu ambil $y = 1 - x$)}.
  \item Turunkan kesamaan ragamnya $\sum_k \bigl(\frac kn -
        x\bigr)^2 p_k(x) = \frac{x(1-x)}{n}$ lalu, lewat
        Cauchy--Schwarz, batas momen pertamanya
        \[
        \sum_{k=0}^{n}\Bigl|\frac kn - x\Bigr|\,p_k(x)
        \leq \sqrt{\frac{x(1-x)}{n}} \leq \frac{1}{2\sqrt n} .
        \]
  \item Tunjukkan bahwa jika $f$ cembung maka $B_nf \geq f$ pada
        $\intcc{0}{1}$ \emph{(lewat ketaksamaan Jensen hingga untuk
        bobot $p_k(x)$)}.
  \item (Batas pencacahan Chebyshev, dinyatakan ulang) Untuk $\delta > 0$
        tunjukkan
        \[
        \sum_{\abs{k/n - x} > \delta} p_k(x)
        \leq \frac{x(1-x)}{n\delta^2}
        \leq \frac{1}{4n\delta^2} ,
        \]
        lalu berikan pembacaan bercorak peluangnya: $B_nf(x)$
        merata-ratakan $f$ atas rata-rata sampel binomial yang memusat di $x$.
\end{enumerate}

\medskip
\noindent\textbf{Bagian II --- Laju: modulus kekontinuan.}
Untuk $\delta > 0$ tetapkan $\omega(\delta) = \sup\{\abs{f(s) - f(t)}
: s, t \in \intcc{0}{1},\ \abs{s - t} \leq \delta\}$.
\begin{enumerate}[resume]
  \item Tunjukkan: $\omega$ hingga, tidak turun,
        $\omega(\delta) \to 0$ ketika $\delta \to 0^+$ (lewat Heine),
        subaditif ($\omega(\delta_1 + \delta_2) \leq
        \omega(\delta_1) + \omega(\delta_2)$), dan
        $\omega(\lambda\delta) \leq (1 +
        \lambda)\,\omega(\delta)$ untuk setiap $\lambda > 0$.
  \item Buktikan taksiran induknya, untuk setiap $\delta > 0$:
        \[
        \abs{B_nf(x) - f(x)}
        \leq \sum_k \omega\Bigl(\Bigl|\frac kn -
        x\Bigr|\Bigr)p_k(x)
        \leq \Bigl(1 + \frac1\delta\sum_k\Bigl|\frac kn -
        x\Bigr|p_k(x)\Bigr)\,\omega(\delta) .
        \]
  \item Ambil $\delta = n^{-1/2}$ lalu simpulkan
        \emph{teorema Weierstrass yang kuantitatif}:
        \[
        \norm{B_nf - f}_\infty \leq
        \frac32\,\omega\Bigl(\frac{1}{\sqrt n}\Bigr)
        \xrightarrow[n\to\infty]{} 0 .
        \]
  \item Turunkan lajunya: $\norm{B_nf - f}_\infty \leq
        \frac{3L}{2\sqrt n}$ bila $L$ konstanta \omterm{def:b2:metric:continuity}{Lipschitz} bagi $f$, dan $\leq
        \frac32 C n^{-\alpha/2}$ bila $\alpha$ eksponen Hölder bagi $f$
        (yakni $\abs{f(s) - f(t)} \leq C\abs{s-t}^\alpha$).
  \item (Contoh yang tajam --- sebuah kesamaan binomial) Untuk $m
        \geq 1$ buktikan
        \[
        \sum_{k=m+1}^{2m} (k - m)\binom{2m}{k}
        = \frac{m}{2}\binom{2m}{m},
        \qquad\text{sehingga}\qquad
        \sum_{k=0}^{2m}\abs{k - m}\binom{2m}{k}
        = m\binom{2m}{m}
        \]
        \emph{(pakailah $k\binom{2m}k = 2m\binom{2m-1}{k-1}$ dan
        kesetangkupan baris binomialnya, yang memberi
        $\sum_{j=m}^{2m-1}\binom{2m-1}{j} = 2^{2m-2}$)}.
  \item Untuk $f(t) = \abs{t - \frac12}$ turunkan nilai eksaknya
        beserta asimtotiknya (lewat binomial pusat,
        \cref{ex:b2:comparison:centralbinomial}):
        \[
        B_{2m}f\Bigl(\frac12\Bigr) - f\Bigl(\frac12\Bigr)
        = \frac{\binom{2m}{m}}{2\cdot4^{m}}
        \;\sim\; \frac{1}{2\sqrt{\pi m}} :
        \]
        jadi laju $\omega(n^{-1/2})$ pada pertanyaan 8 tercapai
        (kecuali sebuah konstanta) --- untuk $f$ yang sekadar \omterm{def:b2:metric:continuity}{kontinu},
        $n^{-1/2}$ milik Bernstein memang jujur.
\end{enumerate}

\medskip
\noindent\textbf{Bagian III --- Teorema Korovkin.} Misalkan $(L_n)$
sebuah barisan operator \emph{linear positif} dari
$C(\intcc{0}{1})$ ke dirinya sendiri sehingga $L_ne_j \to e_j$
secara \omterm{def:b2:funcseq:def}{seragam} untuk $j = 0, 1, 2$.
\begin{enumerate}[resume]
  \item Tunjukkan bahwa $L$ yang linear positif bersifat monoton dan
        memenuhi $\abs{Lf} \leq L\abs f$ \omterm{def:b2:funcseq:def}{titik demi titik}.
  \item Tunjukkan: untuk setiap $\varepsilon > 0$ ada $\delta > 0$
        sehingga untuk \emph{setiap} $s, x \in \intcc{0}{1}$:
        \[
        \abs{f(s) - f(x)} \leq \varepsilon +
        \frac{2M}{\delta^2}(s - x)^2
        \]
        \emph{(tangani $\abs{s - x} \leq \delta$ lewat Heine dan
        $\abs{s-x} > \delta$ lewat batas kasar $2M$)}.
  \item Tetapkan $x$, terapkan $L_n$ pada ketaksamaan pertanyaan 13
        dalam peubah $s$, lalu turunkan
        \[
        \abs{L_nf(x) - f(x)\,L_ne_0(x)}
        \leq \varepsilon\,L_ne_0(x) + \frac{2M}{\delta^2}
        \bigl(L_ne_2(x) - 2x\,L_ne_1(x) + x^2 L_ne_0(x)\bigr).
        \]
  \item Tunjukkan bahwa $\sup_x \bigl(L_ne_2(x) - 2x\,L_ne_1(x) + x^2
        L_ne_0(x)\bigr) \to 0$, lalu rakitlah
        \emph{teorema Korovkin}: yakni $L_nf \to f$ secara \omterm{def:b2:funcseq:def}{seragam} untuk
        \emph{setiap} $f \in C(\intcc{0}{1})$.
  \item Periksalah bahwa $(B_n)$ memenuhi hipotesis Korovkin:
        jadi Weierstrass untuk ketiga kalinya, dari tiga monomial.
  \item Misalkan $I_n$ operator interpolasi afin
        sepotong-sepotong pada simpul $\frac kn$. Tunjukkan bahwa $I_n$
        bersifat linear positif, $I_ne_0 = e_0$, $I_ne_1 = e_1$, dan
        $\norm{I_ne_2 - e_2}_\infty = \frac{1}{4n^2}$
        \emph{(sebab pada tiap sel, galat interpolasi afin
        bagi $t^2$ adalah $(t - a)(b - t)$)}. Lalu simpulkan lewat
        Korovkin: penginterpolasi poligonal konvergen \omterm{def:b2:funcseq:def}{seragam} untuk
        setiap $f$ yang \omterm{def:b2:metric:continuity}{kontinu}.
\end{enumerate}

\medskip
\noindent\textbf{Bagian IV --- Panennya: kepadatan, momen, dan
turunan.}
\begin{enumerate}[resume]
  \item Tunjukkan bahwa polinomial berkoefisien \emph{rasional}
        bersifat padat di $\bigl(C(\intcc{0}{1}),
        \norm\cdot_\infty\bigr)$: jadi \omterm{def:b2:nvs:banach}{ruang Banach} ini
        terpisahkan.
  \item (Momen menentukan fungsinya) Misalkan $f \in
        C(\intcc{0}{1})$ dengan $\int_0^1 f(t)\,t^n \dd t = 0$
        untuk setiap $n \in \N$. Tunjukkan $\int_0^1 f P = 0$ untuk
        setiap polinomial, lalu $\int_0^1 f^2 = 0$, lalu $f = 0$.
  \item Buktikan kesamaan turunannya
        \[
        (B_nf)'(x) = n\sum_{k=0}^{n-1}\Bigl(
        f\Bigl(\frac{k+1}{n}\Bigr) -
        f\Bigl(\frac kn\Bigr)\Bigr)\,
        \binom{n-1}{k}x^k(1-x)^{n-1-k}
        \]
        \emph{(turunkan $p_k$ lalu indeks ulang --- yakni sebuah
        \omterm{thm:b2:series:abel}{penjumlahan Abel})}.
  \item Andaikan $f$ berkelas $C^1$. Dengan memakai teorema nilai
        rata-rata pada tiap pertambahannya lalu membandingkannya dengan
        $B_{n-1}(f')$, tunjukkan bahwa $(B_nf)' \to f'$ secara \omterm{def:b2:funcseq:def}{seragam}
        pada $\intcc{0}{1}$. Turunkan: untuk $f \in C^1$ ada polinomial
        yang konvergen ke $f$ \emph{beserta} turunannya.
  \item Andaikan $f$ berkelas $C^2$. Lewat Taylor--Lagrange di $x$
        tunjukkan
        \[
        \abs{B_nf(x) - f(x)} \leq
        \frac{\norm{f''}_\infty}{2}\cdot\frac{x(1-x)}{n}
        \leq \frac{\norm{f''}_\infty}{8n} :
        \]
        jadi kemulusannya meningkatkan lajunya dari $n^{-1/2}$ menjadi
        $n^{-1}$.
\end{enumerate}

\medskip
\noindent\textbf{Bagian V --- Penjenuhan: teorema
Voronovskaya.}
\begin{enumerate}[resume]
  \item Buktikan kesamaan momen keempatnya
        \[
        \sum_k (k - nx)^4 p_k(x)
        = nx(1-x)\bigl(1 + 3(n-2)x(1-x)\bigr) \leq n^2
        \quad (n \geq 1)
        \]
        \emph{(uraikan $k^4$ atas faktorial turun $k(k-1)\cdots$
        lalu pakai siasat penurunan pertanyaan 2 dua kali
        lagi)}.
  \item (Voronovskaya) Misalkan $f$ berkelas $C^2$ dan $x \in
        \intcc{0}{1}$. Dengan menulis $f(t) = f(x) + f'(x)(t-x) +
        \frac{f''(x)}2(t-x)^2 + \eta(t)(t-x)^2$ dengan $\eta$
        terbatas dan $\eta(t) \to 0$ ketika $t \to x$, buktikan
        \[
        n\bigl(B_nf(x) - f(x)\bigr)
        \xrightarrow[n\to\infty]{}
        \frac{x(1-x)}{2}\,f''(x)
        \]
        \emph{(pecahlah jumlah $\eta$-nya di $\abs{t - x} \leq
        \delta$; lalu kendalikan bagian jauhnya dengan pertanyaan 23)}.
        Jadi galat pertanyaan 22 eksak baik ordenya \emph{maupun}
        konstantanya: $B_n$ \emph{menjenuh} pada $\frac1n$, semulus apa
        pun $f$ --- bandingkan dengan
        \cref{exo:b2:funcseq:7}.
  \item Rangkuman. Satu kalimat untuk masing-masing: (i) apa yang
        dibeli kepositifan semata (Bagian I dan III); (ii) di mana
        kekompakan $\intcc{0}{1}$ masuk pada tiap bagiannya; (iii)
        mengapa tiga fungsi uji sudah cukup pada teorema Korovkin; (iv)
        pertukaran yang dilakukan Bernstein (yakni $n^{-1/2}$ yang
        tangguh untuk $f$ yang kasar, tetapi langit-langit $\frac1n$
        untuk $f$ yang mulus), dan bab mana pada buku ini yang akan
        memainkan permainan sama dengan polinomial trigonometri.

\end{enumerate}
\end{problem}
